\section{Experimental Setup and Metrics}
\label{sec:setup}
\textbf{Hardware and software.} All reported numbers were produced on a laptop with an Intel Core i5-10310U (4 cores, 8 threads), 15~GB RAM and no GPU, with Python 3.10, PyTorch 2.5.1 (CPU), Transformers 4.46.3 and PEFT 0.13.2, all pinned in the repository lock file. Only one large model is loaded at a time (a file lock enforces this). Each run writes a record (config, environment, metrics) under \texttt{results/<area>/}, and the commit SHA is stored with it.

\textbf{Compared systems.} (i) The base Qwen2.5-1.5B-Instruct with the same prompt contract (\emph{before}); (ii) the base plus the shared LoRA adapter, merged into fp32 weights (\emph{after}; merged and unmerged outputs are identical on our check, $|\Delta\text{logit}|\le5.8\times10^{-5}$). Both are evaluated on identical examples with the same parser and the same decoding: plain greedy (\texttt{do\_sample=False}), at most 48 new tokens, fp32, batch 1, with \texttt{repetition\_penalty=1.0}. That last setting matters. Qwen's \texttt{generation\_config.json} ships a repetition penalty of 1.1 that Hugging Face applies \emph{even in greedy mode}; our first validation pass inherited it unnoticed. We kept that pass as an ablation and re-ran everything with the penalty switched off, which also makes the evaluation decoding identical to our hand-written decoders. Invalid or unparsable outputs are counted as an explicit INVALID class and are never dropped. We tuned no hyperparameters, prompts or thresholds: the prompt is fixed by the adapter's training contract, and the full configuration (model and adapter revisions and hashes, decoding settings) was frozen in \texttt{configs/final.yaml} before the first test access. Because of the deadline, the test runs were executed before the final validation run. This leaks nothing, since nothing is selected on validation. \texttt{results/test\_access.log} records every access to the sealed test split.

\textbf{Classification metrics.} We report accuracy, macro-F1, per-class precision/recall/F1 and confusion matrices, overall and separately for inclusion and exclusion criteria, each over its own four-way label set; the overall view uses the union of the six labels. Labels are never collapsed to binary. Macro-F1 is the unweighted mean over the classes of the label set that occur in the gold labels or the predictions. A class absent from both has undefined F1 and is not counted as zero; the fixed-list variant is stored alongside in the metric files. INVALID predictions count as errors.

\textbf{Evidence metrics.} Let $G$ be the expert sentence-ID set and $P$ the predicted set (IDs not present in the note are removed and counted separately as hallucinated). \emph{Micro} precision/recall/F1 pool $|P\cap G|$, $|P|$ and $|G|$ over all records, so records where both sets are empty contribute nothing. \emph{Macro} averages per-record scores with explicit rules for empty sets:
\begin{itemize}
\item $G=P=\emptyset$ scores $P=R=F_1=1$ (a correct abstention);
\item $G=\emptyset$, $P\neq\emptyset$ scores $P=0$, $R=1$, $F_1=0$;
\item $G\neq\emptyset$, $P=\emptyset$ scores $P=1$, $R=0$, $F_1=0$.
\end{itemize}
We also report the macro average over non-empty-gold records only, and the count of records in each of the four empty/non-empty cells.

\textbf{Systems metrics.} Latency per generated token and end-to-end per criterion (median and spread over repeated runs after warm-up), peak resident memory, speed-up of cache vs.\ recomputation, peak memory and time of tiled vs.\ reference attention across sequence lengths, and for speculative decoding the acceptance rate, mean accepted tokens per verify pass and wall-clock speed-up over target-only greedy decoding. Output equality with the baseline is checked token by token.

\textbf{Uncertainty.} With 8 test patients, metric differences are noisy; we report counts alongside percentages and, where available, 95\% percentile bootstrap intervals (1{,}000 resamples of criteria, seed 42).
